\documentclass[10pt,a4paper]{article}
\usepackage{fontspec,amsmath,amssymb,graphicx}
\usepackage[margin=18mm]{geometry}
\setmainfont{Times New Roman}
\setlength{\parindent}{0pt}
\setlength{\parskip}{5pt}
\allowdisplaybreaks
\sloppy
\emergencystretch=2em
\begin{document}
\section*{BT\_4\_6\_3\_04 — Lời giải bổ sung}
\subsection*{Kiến thức và trình tự giải}
\textbf{Bước 1 — Phân biệt đạo hàm cấp hai, ba và bốn.} Quy ước $\widehat f(\xi)=\int e^{-i\xi x}f(x)dx$, biến đổi ngược có hệ số $1/(2\pi)$. Khi đó $\widehat{f^{(m)}}=(i\xi)^m\widehat f$, $\widehat{xf}=i\partial_\xi\widehat f$, $\widehat{xf_x}=-\widehat f-\xi\partial_\xi\widehat f$.

\textbf{Bước 2 — Giải phương trình của biến đổi.} $u_t=u_{xx}$ cho $\widehat u=e^{-t\xi^2}\widehat f$. $u_t=u_{xxx}$ cho $\widehat u=e^{-it\xi^3}\widehat f$. $u_t=-u_{xxxx}$ cho $\widehat u=e^{-t\xi^4}\widehat f$. Với $xu_x$, giải PDE cấp một theo $\xi$ bằng đường đặc trưng, gồm cả hệ số biên độ $-\widehat u$.

\textbf{Bước 3 — Biến đổi ngược và ý nghĩa tích phân.} Nhân bậc ba là $A_t(x)=(3t)^{-1/3}\operatorname{Ai}(-x/(3t)^{1/3})$, tích phân dao động. Nhân bậc bốn là $B_t(x)=(2\pi)^{-1}\int e^{i\xi x-t\xi^4}d\xi$, hội tụ tuyệt đối theo $\xi$ và có thể đổi dấu. Không áp dụng tính dương của Gauss hay nguyên lý cực đại cấp hai cho hai nhân này.

\textbf{Bước 4 — Miền hữu hạn cần kiểm tra biên và toán tử liên hợp.} Với $u_t=Ku_{xxx}$, $u(0)=u_x(0)=u(L)=0$, năng lượng cho $I'=-K u_x(L)^2/2$. Toán tử liên hợp là $-KD^3$ với $Y(0)=Y(L)=Y'(L)=0$. Đây không phải bài Sturm–Liouville tự liên hợp cấp hai; không suy mọi trị riêng đều thực và không tự dùng hệ số $\int fX/\int X^2$. Có thể giải bằng biến đổi Laplace theo thời gian và nghiệm cơ bản của ODE bậc ba, tránh khẳng định cơ sở chưa chứng minh.

\textbf{Bước 5 — Nêu lớp dữ liệu đủ cho công thức.} Với bậc ba trên toàn trục, công thức Fourier hợp lệ trước hết cho dữ liệu Schwartz hoặc $L^2$ theo nghĩa biến đổi Fourier. Dữ liệu liên tục bị chặn là phân phối tempered; phép nhân $e^{-it\xi^3}\widehat f$ vẫn hợp lệ vì các đạo hàm của thừa số tăng đa thức, nên xác định nghiệm trong $\mathcal S\prime$. Tuy nhiên điều đó không bảo đảm tích phân Airy hội tụ tuyệt đối hay nghiệm tự động là hàm cổ điển. Ghi đúng ý nghĩa của công thức theo lớp dữ liệu.
\subsection*{Lời giải áp dụng cho bài}
\textbf{Bước 1. Quy ước Fourier và các công thức dùng chung.}
$\widehat f(\xi)=\int e^{-i\xi x}f(x)dx$, biến đổi ngược $(2\pi)^{-1}\int e^{i\xi x}\widehat f(\xi)d\xi$. Tích phân từng phần cho $\widehat{u_{xx}}=-\xi^2\widehat u$, $\widehat{u_{xxx}}=-i\xi^3\widehat u$, $\widehat{u_{xxxx}}=\xi^4\widehat u$. Đạo hàm theo $\xi$ cho $\widehat{xu_x}=-\widehat u-\xi\widehat u_\xi$. Ký hiệu $g_K(z,t)=(4\pi Kt)^{-1/2}e^{-z^2/(4Kt)}$.

\textbf{(i) $u_t=u_{xx}+xu_x$.}
Đặt $W=\widehat u$: $W_t+\xi W_\xi=-(1+\xi^2)W$. Đặc trưng $\xi(s)=\xi_0e^s$ cho
\[W(\xi_0e^t,t)=e^{-t}\exp[-\xi_0^2(e^{2t}-1)/2]\widehat f(\xi_0).\]
Vì $\xi_0=e^{-t}\xi$,
\[W=e^{-t}e^{-\xi^2(1-e^{-2t})/2}\widehat f(e^{-t}\xi).\]
Đổi biến $\eta=e^{-t}\xi$ trong biến đổi ngược, đặt $z=e^tx,\tau=(e^{2t}-1)/2$, được
\[u(x,t)=(g_1(\cdot,\tau)*f)(e^tx).\]
Kiểm tra: $u_t=zv_z+e^{2t}v_\tau$, $u_{xx}=e^{2t}v_{zz}$, $xu_x=zv_z$, nên $v_\tau=v_{zz}$.

\textbf{(ii) $u_t=u_{xxx}$.}
$W_t=-i\xi^3W$ cho $W=e^{-it\xi^3}\widehat f$. Do đó
\[u=\mathcal F^{-1}(e^{-it\xi^3}\widehat f)=A_t*f,\qquad
A_t(x)=(3t)^{-1/3}\operatorname{Ai}(-x/(3t)^{1/3}).\]
Để thấy dấu Airy, đổi $\eta=(3t)^{1/3}\xi$ trong $(2\pi)^{-1}\int e^{i\xi x-it\xi^3}d\xi$ rồi đổi $\eta\mapsto-\eta$, được tích phân Airy có pha $\eta^3/3-\eta x/(3t)^{1/3}$. Với dữ liệu Schwartz, công thức cho nghiệm cổ điển; với $L^2$ là nghiệm Fourier trong $L^2$. Với dữ liệu chỉ liên tục bị chặn như đề, $f$ là phân phối tempered và phép nhân $e^{-it\xi^3}\widehat f$ hợp lệ vì các đạo hàm của thừa số có tăng trưởng đa thức; kết quả được hiểu trong $\mathcal S'$. Không khẳng định tích phân $A_t*f$ hội tụ tuyệt đối hoặc nghiệm tự động là hàm cổ điển.

\textbf{(iii) $u_t=-u_{xxxx}+xu_x$.}
$W_t+\xi W_\xi=-(1+\xi^4)W$. Cùng đặc trưng $\xi=\xi_0e^t$ cho
\[W=e^{-t}e^{-\xi^4(1-e^{-4t})/4}\widehat f(e^{-t}\xi).\]
Đặt $z=e^tx$, $\tau=(e^{4t}-1)/4$ và
\[B_\tau(z)=\frac1{2\pi}\int_{\mathbb R}e^{i\xi z-\tau\xi^4}d\xi.\]
Đổi biến Fourier ngược cho $u=(B_\tau*f)(e^tx)$. Kiểm tra $u_t=zv_z+e^{4t}v_\tau$, $u_{xxxx}=e^{4t}v_{zzzz}$, nên $v_\tau=-v_{zzzz}$. $B_\tau$ là hàm Schwartz với $\tau>0$, nhưng không phải nhân Gauss dương.

\textbf{(iv) $u_t=u_{xx}-xu_x$.}
$W_t-\xi W_\xi=(1-\xi^2)W$. Đặc trưng $\xi=\xi_0e^{-t}$ cho $W=e^t e^{-\xi^2(e^{2t}-1)/2}\widehat f(e^t\xi)$. Đổi biến trong Fourier ngược, đặt $z=e^{-t}x$, $\tau=(1-e^{-2t})/2$:
\[u(x,t)=(g_1(\cdot,\tau)*f)(e^{-t}x).\]
Kiểm tra $u_t=-zv_z+e^{-2t}v_\tau$, $u_{xx}=e^{-2t}v_{zz}$, $-xu_x=-zv_z$, nên đúng PDE. Ở (i),(iii),(iv), khi $t\downarrow0$, $z\to x,\tau\to0$, các nhân trả lại $f$.

\textbf{(v) $u_t=u_{xx}+tx$ trên nửa trục.}
Bản in lặp lại $u(x,0)=0$ hai lần, không cho biên tại $x=0$. Một nghiệm riêng là $u_p=xt^2/2$: $(u_p)_t=tx$, $(u_p)_{xx}=0$, $u_p(x,0)=0$. Nếu bổ sung $u(0,t)=0$ và lớp tăng trưởng phù hợp, đây là nghiệm của bài bổ sung. Với bản in, còn họ nghiệm
\[u=\frac{xt^2}2+\int_0^t
\frac{x e^{-x^2/[4(t-s)]}}{2\sqrt\pi(t-s)^{3/2}}\mu(s)ds,\]
với $\mu$ trơn tùy ý vanishing gần $0$. Hạng thêm giải nhiệt, dữ liệu đầu $0$ và có biên $\mu$, nên chứng minh không duy nhất khi thiếu biên. Dòng rời về $f=\chi_{(0,1)}$ không thống nhất với dữ liệu đầu $0$ đã in; không tự đưa nó vào nghiệm của hệ này.

\textbf{(vi) Dirichlet nửa trục, hệ số $k>0$.}
Thác triển lẻ $f(y)=1/(1+y^2)$ cho
\[u=\int_0^\infty
\frac{e^{-(x-y)^2/(4kt)}-e^{-(x+y)^2/(4kt)}}{\sqrt{4\pi kt}}
\frac{dy}{1+y^2}.\]
Hệ số $k$ phải có trong cả mẫu và số mũ. Hai ảnh triệt tiêu tại $x=0$; Gauss xấp xỉ đơn vị trả lại $f$ khi $x>0,t\downarrow0$; đạo hàm nhân cho $u_t=ku_{xx}$. Giá trị đầu $f(0)=1$ không khớp biên $0$ tại góc.

\textbf{(vii),(viii) Nguồn trên nửa trục nhưng thiếu biên.}
Đặt $H(x,y,t)=g_1(x-y,t)-g_1(x+y,t)$ và $P(x,t)=xe^{-x^2/(4t)}/(2\sqrt\pi\,t^{3/2})$. Với một biên Dirichlet được \emph{bổ sung} là $\mu(t)$, công thức Duhamel cho
\[u=\int_0^\infty H(x,y,t)f(y)dy+
\int_0^t\!\int_0^\infty H(x,y,t-s)F(y)dy\,ds+
\int_0^tP(x,t-s)\mu(s)ds.\]
Ở (vii), $f(y)=1/(1+y^2)$, $F(y)=e^{-y}$; ở (viii), $f=0$, $F(y)=e^{-y}/(1+y^2)$. Để kiểm tra hạng biên, đổi $r=t-s$:
\[\int_0^tP(x,r)dr=\operatorname{erfc}\frac{x}{2\sqrt t}\longrightarrow1
\quad(x\downarrow0).\]
Khối lượng của $P(x,r)$ tập trung tại $r=0$ khi $x\downarrow0$, nên hạng này tiến về $\mu(t)$ nếu $\mu$ liên tục. Nó giải nhiệt và có mặt đầu $0$ với $x>0$. Chọn hai hàm $\mu$ khác nhau vanishing gần $0$ cho hai nghiệm cùng PDE và dữ liệu đầu, khác biên; do đó đề gốc không xác định duy nhất. Các tích phân không gian với $1/(1+y^2)$ là biểu diễn chính xác; chúng không phải tích phân sơ cấp cần giả tạo một nguyên hàm.

\textbf{(ix) Neumann, dữ liệu $e^{-x}$, $K=1$.}
Nhân ảnh chẵn cho $u=\int_0^\infty[g_1(x-y,t)+g_1(x+y,t)]e^{-y}dy$. Hoàn thành bình phương:
\[-(x\mp y)^2/(4t)-y=-(y+2t\mp x)^2/(4t)+t\mp x.\]
Đổi biến $z=(y+2t\mp x)/\sqrt{4t}$ rồi dùng erfc:
\[u=\frac12e^{t-x}\operatorname{erfc}\frac{2t-x}{\sqrt{4t}}
+\frac12e^{t+x}\operatorname{erfc}\frac{2t+x}{\sqrt{4t}}.\]
Hàm chẵn theo $x$ cho $u_x(0,t)=0$, hai giới hạn erfc cho $u(x,0)=e^{-x}$ với $x>0$; nhân ảnh kiểm tra PDE.

\textbf{(x) Dirichlet bằng $1$.}
Đặt $w=u-1$. Khi đó $w_t=w_{xx}$, $w(0,t)=0$, $w(x,0)=1/(1+x^2)-1$. Thác triển lẻ cho
\[u=1+\int_0^\infty H(x,y,t)\left(\frac1{1+y^2}-1\right)dy.\]
Tại $x=0$ nhân $H$ bằng $0$, nên $u=1$. Mặt đầu trả lại $f$, và phép cộng hằng không đổi PDE. Dữ liệu góc về giá trị khớp vì $f(0)=1$.\end{document}